\documentclass[tikz,border=8pt]{standalone}
\usepackage{tikz}
\usepackage{amsmath}
\usepackage{pgfplots}
\pgfplotsset{compat=1.18}
\usepgfplotslibrary{groupplots}
\usepackage{ctex}
\begin{document}
\begin{tikzpicture}
\begin{groupplot}[
  group style={group size=2 by 1, horizontal sep=1.5cm},
  xlabel style={font=\small},
  ylabel style={font=\small},
  xticklabel style={font=\scriptsize},
  yticklabel style={font=\scriptsize},
  axis lines=left,
  samples=100,
  clip=false,
]

% Left: PDF comparison
\nextgroupplot[
  width=6.5cm, height=5.5cm,
  xmin=0, xmax=4.5,
  ymin=0, ymax=2.2,
  xlabel={$x$},
  ylabel={$f(x)$},
  title={指数分布 PDF：$f(x)=\lambda e^{-\lambda x}$},
  title style={font=\small\bfseries},
  legend pos=north east,
  legend style={font=\scriptsize},
]
\addplot[blue, very thick, domain=0:4.5] {0.5*exp(-0.5*x)};
\addlegendentry{$\lambda=0.5$，$E=2$}
\addplot[red!80, very thick, domain=0:4.5] {exp(-x)};
\addlegendentry{$\lambda=1$，$E=1$}
\addplot[green!60!black, very thick, domain=0:4.5] {2*exp(-2*x)};
\addlegendentry{$\lambda=2$，$E=0.5$}

% Right: memoryless illustration
\nextgroupplot[
  width=6.5cm, height=5.5cm,
  xmin=0, xmax=5,
  ymin=0, ymax=1.05,
  xlabel={$t$},
  ylabel={$P(X > t)$},
  title={无记忆性：生存函数 $P(X>t)=e^{-\lambda t}$},
  title style={font=\small\bfseries},
  legend pos=north east,
  legend style={font=\scriptsize},
]
\addplot[blue, very thick, domain=0:5] {exp(-0.5*x)};
\addlegendentry{$\lambda=0.5$}
\addplot[red!80, very thick, domain=0:5] {exp(-x)};
\addlegendentry{$\lambda=1$}
\addplot[green!60!black, very thick, domain=0:5] {exp(-2*x)};
\addlegendentry{$\lambda=2$}
% Annotate memoryless: vertical dashed at s=1
\draw[dashed, gray] (axis cs:1,0) -- (axis cs:1,1.0);
\node[font=\scriptsize, fill=white, inner sep=1pt] at (axis cs:1,1.02) {$s$};
\node[font=\scriptsize, text=red!80, text width=3cm, align=center]
  at (axis cs:3.5, 0.75)
  {$P(X{>}s{+}t\,\vert\, X{>}s)$\\$=P(X{>}t)$};

\end{groupplot}
\node[font=\scriptsize, fill=yellow!15, draw, rounded corners=2pt, text width=8cm]
  at (7,-0.9) {指数分布是连续分布中唯一具有无记忆性的分布（类比离散的几何分布）};
\end{tikzpicture}
\end{document}
